\section{降低张量幂的总元数}\label{sec:rounds}

引理~\ref{lem:local} 的集中估计把网络分为少数大站点和一个由小站点组成的核心。要降低网络的求值成本，还必须降低\emph{总}元数。事实上，若站点 $v$ 的元数为 $u_v$，则引理~\ref{lem:batch-table} 的批量求值分摊到每个张量副本上的成本至多是站点数与下式之和的常数倍：\[
       \sum_{v:\,u_v>0}u_v^p
           \le \bigl(\max_v u_v\bigr)^{p-1}\sum_v u_v .
\] 这里 $1<p<2$。因此，只要能消去足够多的关联比特，就允许最大元数有所增长。下面的构造反复利用这一权衡，同时控制元数的平方和，使每个新网络仍满足集中估计的假设。

本部分始终设 $k$ 是满足 $\omega_k<9/4$ 的有限域，所有对数均以二为底。域 $K$ 上的\emph{表网络}以比特为连线，每个元数为 $s$ 的站点配置一张表 $\{0,1\}^{s}\to K$。网络的规模是站点数加上各站点的元数之和，外腿也计入元数。两个站点之间并行的比特组成一个有序的\emph{线束}。为网络固定一部字典，每个站点的标签指明其表在字典中的名称。

\needspace{12\baselineskip} \begin{lemma}[张量幂的压缩]\label{lem:macro} 设 $h=\lceil\log_2 n\rceil$，$0<\delta\le 1/10$，且 $0<\beta\le b\le 1/2$。对每个仅依赖于 $k,\delta,\beta$ 的充分小的固定 $a>0$，以及一切充分大的 $n$，以下命题成立。取次数为 $O(\log h)$ 的扩域 $K/k$，使其大小足以在下文所用尺度上应用引理~\ref{lem:local}。假设输入表网络的最大元数至多为 $h^b$，表标签的长度为 $(\log h)^{O(1)}$ 比特，规模至多为 $m$，其中 $n\le m\le n^{1.1}$，且外部比特数满足 $e\le m h^{-2a}$。则存在满足 $1\le d=n^{o(1)}$ 的 $d$，使该网络的 $d$ 次张量幂可以表示为 $K$ 上的表网络。其 $de$ 个外部比特均可分别访问，并且 \[
 \begin{aligned}
 \text{规模}&\le md\,h^{-(1-\delta)a},&
 \text{最大元数}&\le h^{b+(1+\delta)a}.
 \end{aligned}
\] 输出标签的长度仍为 $(\log h)^{O(1)}$ 比特。对 $n$ 的下限要求可以依赖于输入标签长度上界中的固定指数。\end{lemma}

称引理~\ref{lem:macro} 的一次应用为一次\emph{宏操作}。操作前，每份副本的查表开销至多为 $O(mh^{b(p-1)})$。宏操作后，该上界乘上因子 \[
       h^{-a\{(1-\delta)-(1+\delta)(p-1)\}} .
\] 由于 $p<2$，只要 $\delta$ 取得足够小，指数便为负。第~\ref{sec:shaving} 节在对表求值之前先累积这些节省。

\begin{table}[bp]\centering \small \begin{tabular}{@{}ll@{}} \toprule 量 & 在一个求和项中的含义\\ \midrule $H,\ N_b$ & 公共宿主图的位置数与使用该图的块数\\ $r$ & 每次分组时合并的块数\\ $S,\ T_0$ & 初始块数与每块的原始副本数，满足 $d=ST_0$\\ $L,\ T$ & 总元数预算与当前的元数尺度\\ $D$ & 各轮均保持固定的集中参数\\[2pt] $X,\ Z$ & 核心站点之间的宽度与冻结站点的总计费量\\ $\Delta,\ \nabla$ & 核心宿主图与整个宿主图的度数上界\\ $B,\ u$ & 最大装包负载与内部宽度单位\\ \bottomrule \end{tabular} \caption{压缩张量幂所用的参数。预算 $L$、$X$ 和 $Z$ 各自涵盖该求和项的所有块。}\label{tab:macro} \end{table}

宏操作将原张量的一个幂展开为若干乘积的和。每个乘积包含许多因子，称为\emph{块}。把块按每组 $r$ 个合并，以便把站点装包或重新排列其连线。每一轮先使元数的平方集中，再进行固定次数的连线缩减以缩短核心，最后恢复有界度数，使下一轮可以开始。一次连线缩减称为一\emph{拍}，由若干次收缩和随后的一次核心度数重置组成。轮数随 $h$ 增长。一旦固定 $D$，每轮\emph{内部}的重复次数就都是常数。表~\ref{tab:macro} 列出了各项预算及其衡量的量。

不同的求和项可采用不同的接线方式。第~\ref{sec:rounds} 至第\ref{sec:full} 节中的界适用于每一个求和项。第~\ref{sec:assembly} 节用一个固定网络表示这些项之和，从而完成引理~\ref{lem:macro} 的证明。

\subsection{公共图及其初始负载} 每个求和项都是若干\emph{块}的乘积，各块的外部比特列表互不相交。同一求和项的 $N_b$ 个块共用一张描述可能邻接关系的简单图，称为\emph{宿主图}。宿主图的一个位置和一个块组成一个位置实例。实例可以为空。每条实际的内部连线连接同一块中的两个站点，局部端口顺序指明两端哪些比特相互对应。因此，同一条宿主边在不同块中可以承载不同的线束。允许宿主图依赖于求和项。

每个分组步骤都给块重新编号，每 $r$ 个块分为一组，其中 $r$ 是二的幂。每组内，先前的块相乘，并一起接受变换。新位置可以用一个旧位置、$r$ 个指标的一个二进子区间以及有界个其他选择来标记。这样的区间少于 $2r$ 个。若 $H$ 为宿主图的位置数，则任一步允许填充后的规模满足 \begin{equation}\label{eq:round-occurrence}
 H'\le c rH, \qquad N_b'=N_b/r,
\end{equation} 其中 $c$ 在这一步的各次重复中保持固定。此处的常数可以依赖于单轮方案。在单个位置使用固定数量槽位的变换可以并入相邻的分组步骤。

每轮开始时，每个非空的活动实例包含一个站点。宿主图的度数有绝对常数上界，外部比特固定在指定端口。设 $T$ 为二的幂次宽度，$L$ 为总元数预算。若活动站点的实际元数为 $s_v$，则其计费元数为 $z_v=\max(s_v,\theta T)$，其中 $\theta>0$ 是下文选定的一个小的二进常数。我们保持下列不变量：\begin{equation}\label{eq:round-invariants}
 \sum_v z_v\le L,\qquad \sum_v z_v^2\le LT,\qquad
 \max_v z_v\le DT/2 .
\end{equation} 这三个界都涵盖该求和项的所有块。第一个界用于支付保留下来的比特，第二个界是下一次局部集中的输入条件，第三个界保证把若干站点装包时不会超出单个单元的容量。下限 $\theta T$ 也限制了给定计费量的一包中能纳入多少份表描述。移除所有标量站点，将各自的表项保留为该求和项的因子。

先将输入替换成一张度数有界的图。对每个正元数为 $s$ 的站点，在其端口上设置一棵平衡二叉树。根将 $s$ 个比特作为一个线束接收。每个分叉处的恒等张量把一个有序线束分为两个，长度之差至多为一。在每个叶节点放置恒等张量，并将原有连线及外部比特接在叶上。把所得的两个叶节点连接起来也能处理原输入中的自环。每条树边两端的元数之比有界。叶与叶之间的连线两端宽度均为一。树有 $O(\log(2+h))$ 层，总元数为 $O(m\log(2+h))$。所有新表都有短标签。

选取二的幂 $T_0\in[4h^b,8h^b]$，叠放刚构造的网络的 $T_0$ 个副本。对元数为 $s$ 的站点，令 $v$ 为不超过 $T_0/s$ 的最大二的幂。把副本划为长度 $v$ 的连续区间。每个区间独占一个位置，承载相应各表的乘积。其元数 $sv$ 属于 $(T_0/2,T_0]$。由于每条旧边两端的元数相当，每个区间只与另一端有界多个区间相交。每条边先继承原线束的顺序，再继承副本的顺序，因而两端的比特顺序相同。叠放后的图度数有界，位置数为 $H_0=O(m\log(2+h))$，因为元数为 $s$ 的站点至多产生 $2s$ 个区间。一个区间的表由旧表标签、区间长度和有界多个有序端口长度指定。

取该叠层的 $S$ 个块，将 $H_0$ 填充至 $m\log(2+h)$ 量级。把先前移除的输入标量之积的 $d$ 次幂作为最终张量的因子。初始参数为 \begin{equation}\label{eq:round-start}
 d=T_0S,\qquad T=T_0,\qquad L=L_0=CmT_0S\log(2+h).
\end{equation} 当绝对常数 $C$ 足够大时，总计费元数至多为 $L_0$。每个站点的元数属于 $(T_0/2,T_0]$，故其平方至多为元数的 $T_0$ 倍。这就得到 \eqref{eq:round-invariants} 中的三个界。整个构造过程中，$ed$ 个外部比特各自保留原来的名称。

\section{降低未冻结站点之间的宽度}\label{subsec:core-width}

集中操作之后，有些站点在本轮余下的过程中保持冻结，其余站点构成\emph{核心}。我们先降低核心站点间的宽度，再恢复有界度数，并把所有站点交给下一轮。每个核心位置实例承载一个张量。每个冻结位置实例持有一个张量列表，并记录每个张量冻结时的计费量。所有外腿都接在冻结张量上。冻结站点间的宿主邻接关系不承载任何连线，可以删除。

设 \(X\) 为所有块中核心站点与核心站点间连线的宽度之和，每条线只计一次。设 \(Z\) 为冻结站点的总计费量。核心与冻结站点之间的连线称为\emph{系线}，其总宽度至多为 \(Z\)。装包时，将每个核心位置实例的负载取为其元数，将每个冻结位置实例的负载取为其计费量之和乘以属于 \((0,1]\) 的固定二进因子。用 \(B\) 表示单个负载的上界。所有负载之和满足 \(\mathcal M\le 2X+2Z\)。存在一个共同的整数 \(u\le B\)，使所有内部宽度都是 \(u\) 的倍数。

宿主图记录所有块中可能出现的接触。用 \(\Delta\ge3\) 表示其核心度数的上界，用 \(\nabla\ge\Delta\) 表示整个宿主图的最大度数的上界。核心站点在不同块中可能与不同的冻结站点接触，因此即使核心的度数为三，仍可能有许多潜在系线。下面的宽度估计与潜在系线的条数无关。

\begin{lemma}[核心宽度的缩减]\label{lem:core-width} 固定输入核心度数的一个上界，并设定理~\ref{thm:sort} 的假设成立。存在一个固定的分组及收缩步骤序列，称为一\emph{拍}。给定余量 \(\sigma>0\)，每个求和项都满足 \[
 X_{\rm out}\le \tfrac12 X+C(Z+\sigma),\qquad
 \Delta_{\rm out}\le3,\qquad Z_{\rm out}=Z.
\] 这一恒等式分别保留每条外腿。最大负载、线束单位的倒数和 \(\log(2\nabla)\) 的增长倍数均有固定上界。这些因子及 \(C\) 均与输入的全图最大度数无关。

在宿主图有 \(H\) 个位置、块数 \(N_b\) 为二的幂的一步中，选取二的幂 \(r\mid N_b\)。分组后剩下 \(N_b/r\) 个块，至多使用 \(crH\) 个位置。只须在每步取 \(r\ge c_*\)、\(c_*HN_bB/r\le\sigma\)，并使初始单位能被各细分因子的乘积整除。位置常数及余量常数 \(c,c_*\) 可以依赖于相应步骤的度数上界。\end{lemma}

一拍由若干节省连线的子轮和一次重置组成。在每个子轮中，允许可能的核心接触增多，以换取核心宽度的缩减。重置时，将余下的连线经过一个度数为三的核心布设。我们将选择子轮数，使各子轮的节省合起来足以抵付这次布线的开销。

\subsection{一个接线恒等式} 将六方按图~\ref{fig:two-layouts} 排成三角阵列。若两方同在一行或一列，则称这一对为\emph{允许对}。所需的相等连线 \(AB\) 和 \(CD\) 的两端既不同行也不同列。我们将其乘积表示为若干只在允许对上布线的网络之和，随后按行或按列收缩整个划分。这样，同一个求和项有两种布局可选。

\begin{figure}[htb]\centering \begin{tikzpicture}[x=1.06cm,y=1.32cm,font=\small,
 party/.style={circle,inner sep=1pt,minimum size=4mm},
 part/.style={rounded corners=2.8pt,line width=.65pt},
 demand/.style={draw=ink,densely dashed,thick}]
\foreach \panel/\shift/\heading in {0/0/需求连线,1/4.2/按行划分,2/8.4/按列划分}{
 \begin{scope}[xshift=\shift cm]
 \node[font=\sffamily\small] at (1.14,1.96) {\heading};
 \ifnum\panel=1
  \draw[part,draw=sea,fill=sea!8] (-.33,1.15) rectangle (2.61,1.65);
  \draw[part,draw=sea,fill=sea!8] (-.33,.45) rectangle (1.47,.95);
  \draw[part,draw=sea,fill=sea!8] (-.33,-.25) rectangle (.33,.25);
 \fi
 \ifnum\panel=2
  \draw[part,draw=rust,fill=rust!8] (-.33,-.25) rectangle (.33,1.65);
  \draw[part,draw=rust,fill=rust!8] (.81,.45) rectangle (1.47,1.65);
  \draw[part,draw=rust,fill=rust!8] (1.95,1.15) rectangle (2.61,1.65);
 \fi
 \ifnum\panel=0
  \draw[demand,shorten <=2.6mm,shorten >=2.6mm] (0,1.4)--(1.14,.7);
  \draw[demand,shorten <=2.6mm,shorten >=2.6mm] (2.28,1.4)
    .. controls (2.64,-.14) and (1.26,-.23) .. (0,0);
 \fi
 \foreach \x/\y/\name in {0/1.4/C,1.14/1.4/M,2.28/1.4/A,0/.7/N',1.14/.7/D,0/0/B}{
  \node[party] at (\x,\y) {\(\name\)};
 }
 \end{scope}}
\end{tikzpicture} \caption{例如，\(C\) 和 \(M\) 之间的一条相等连线在按行收缩时消失，在按列收缩时保留。每条允许连线在其中一种布局下都会消失。对于一个求和项中该恒等式的所有副本，统一采用两种布局中代价较低的一种，所保留的宽度至多为替代连线总宽度的一半。}\label{fig:two-layouts} \end{figure}

一条 \(s\) 比特连线是两个 \(s\) 比特指标上的相等张量。对于 \(k\) 上一个尺度为一的普通价格，令 \(m\) 表示允许对上最大的价格距离。在允许图的一棵生成树上，令每个顶点处的关联比特相等，就能提供一个六方共享的比特。因此 \(1\le5m\)。两个所需距离都至多为 \(2m\)。若 \(d(C,D)=2m\)，则路径 \(C,M,D\) 和 \(C,N',D\) 使 \(M\) 和 \(N'\) 在一个 CAT(0) 扩张中成为同一个中点。若还满足 \(d(A,B)=2m\)，那么 \(C\) 与 \(M=N'\) 都是 \(A,B\) 的中点。于是 \(C=M\)，这与 \(d(C,M)=m>0\) 矛盾。由尺度为一的价格的紧致性，存在统一的 \(\lambda_*<2\)，满足 \[
 d(A,B)+d(C,D)\le2\lambda_*m .
\]

取 \(\lambda_*<\lambda_0<2\)。我们断言，对固定的 \(\eta'>0\) 和任意大的 \(b\)，两条宽度为 \(b\) 的需求连线可以写成至多 \(2^{\eta'b}\) 个网络的普通和，每个网络的允许连线总宽度至多为 \(2\lambda_0b\)。事实上，选充分大的 \(h_0\) 使 \(\lfloor2\lambda_0h_0\rfloor>2\lambda_*h_0\)。作为供给，在每个允许对上取一条宽度为 \(\lfloor2\lambda_0h_0\rfloor\) 的连线，再对所有允许对取直和。在每个谱点，这一供给的价格都超过两条宽度为 \(h_0\) 的需求连线的价格。由严格谱比较，对 \(b=h_0r\)，宽度为 \(b\) 的需求是供给的 \(r\) 次幂的有界多个副本之直和的一个限制。展开直和并施以局部限制映射，即得所称网络。记允许对的集合为 \(E_{\rm perm}\)，一个求和项的指标需要 \(r\log_2|E_{\rm perm}|+O_{h_0}(1)\) 比特。先取充分大的 \(h_0\)，再取充分大的 \(r\)，使上述比特数不超过 \(\eta'b\)。\(k\) 上的该恒等式可扩张至任意扩域 \(K/k\)。为构造子轮，我们需要即使每条线承载同一比特的许多副本，选择的数目仍然固定。

\begin{lemma}[宽度的固定选择]\label{lem:powered-gadget} 存在常数 $b\in\mathbb N$ 和 $\lambda<2$，使得对每个 $z\ge1$，两条宽度为 $bz$ 的需求连线等于同一组六方上有界多个网络的和。每个网络在允许边上至多使用 $2\lambda b$ 个线束，每束含 $z$ 比特。局部张量由 $z$、所在的一方及求和项指标指定。指标集及各宽度的选取（以 $z$ 为单位）可以独立于 $z$ 固定。\end{lemma}

\begin{proof} 将上述有限恒等式取 $z$ 次张量幂。起初，每一项由含 $z$ 个选择符的序列标记，每个选择符至多占 $\lceil\eta'b\rceil$ 比特。在每条允许边上，把比特总数向上取整至 $z$ 的倍数。按取整后以 $z$ 为单位的各边宽度向量对序列分组，允许空组。该向量取值于固定有限集，每条允许边为取整多付至多一个单位。每一组中，取连通的允许图的一棵生成树，沿每条树边传送完整序列，每边额外使用至多 $\lceil\eta'b\rceil$ 个单位。局部表检验所有关联树边上的序列是否相同且有效、是否属于该组，并检验填充位置是否为零。它使用该序列将每条原边的位置划为规定的片段，再计算 $z$ 张局部表的乘积。由连通性，一旦指定序列，它在各树边上的编码就是唯一的。因此，对树变量求和恰好给出该组的和。以 $z$ 为单位的总数至多为 $2\lambda_0 b+|E_{\rm perm}|+5\lceil\eta'b\rceil$。先取充分小的 $\eta'$，再取充分大的 $b$，就能使该量至多为 $2\lambda b$，其中 $\lambda<2$。\end{proof}

\subsection{子轮} 在一个子轮中，各个位置的位置实例被分别装包。来自不同包的两个出向单位可以分别充当 \(AB\) 和 \(CD\)。先安排足够多这样的配对，并保证单次收缩所合并的包数有界。

\paragraph{平衡划分。} 对每一组 \(r\) 个块指标，按共同顺序取二进区间。给每条宿主边 \(e\) 指定传播距离 \[
 s_e=64\lceil\log_2(2x_e)\rceil,\qquad
 x_e=\begin{cases}\Delta&\text{在核心边上},\\
                     \nabla&\text{在系线上}.\end{cases}
\] 在每个位置标记负载大于 \(AB\) 的所有区间，其中 \(2\le A\le r\) 为二的幂。标记意味着必须分裂。将每个标记传播给同一位置的父区间，并且沿每条关联边 \(e\)，传播给另一端高 \(s_e\) 层的祖先区间，只要该祖先存在。叶区间称为聚落，其负载至多为 \(AB\)，长度至少为 \(A\)。由于较小聚落的父区间所发出的传播，边 \(e\) 两端相交的聚落长度之比至多为 \(2^{s_e}\)。

在每组内部对指标采用一个公共置换，可使聚落总数满足 \begin{equation}
 Q\le C_1\left( HN_b/r+\mathcal M/(AB)\right)
 \label{eq:core-colonies}
\end{equation} 其中 \(C_1\) 为绝对常数。为看出这一点，在一个二进制基下把指标等同于 \(\mathbb F_r\)，再乘以均匀随机选取的非零元。两个相异指标落入同一深度 \(j\) 区间的概率至多为 \(2^{-j}\)。对一组中的一个源位置，记 \(x_i\le1/A\) 为各负载除以 \(AB\) 之后的值，令 \(x=\sum_i x_i\)。初始即带标记的区间内，不同归一化负载的有序交叉乘积之和至少为 \(1/2\)。因此，深度 \(j\) 处的初始标记数的期望至多为 \(2x^2 2^{-j}\)。在每个顶点，对关联边求 \(2^{-s_e/16}\) 之和至多为 \(1/2\)。对所得路径权重求和可知，总传播距离至多为 \(s\) 时，可达位置至多有 \(2\cdot2^{s/16}\) 个。固定总向上位移 \(s\)，允许在同一位置向父区间迈步。那么在每个可达位置，深度 \(j\) 处的源标记到达的不同祖先至多有 \(2^{j-s}\) 个。在该位置对 \(j\ge s\) 求和，期望数至多为 \[
 \sum_{j\ge s}\min\{2^{j-s},2x^2 2^{-j}\}=O(x2^{-s/2}).
\] 再对距离、源和组求和，标记总数的期望至多为 \(O(\mathcal M/(AB))\)。每个二叉划分的叶数比带标记的内部节点数多一。选择使实际数目最小的乘子即可。

\paragraph{节省连线的一次子轮。} 为所有核心边定向。用 \(\mathbb F_2\) 上的一个仿射泛函作用于块在组内的指标，将每个块着为粗或细两色之一。同一块的所有位置颜色相同。在同一个位置装包时，取各张量的张量积，同时保留连线。随后合并产生的自环收缩到局部表内。在每个核心聚落，把所有粗张量装在一个包内。其余依次装成细包，负载达到 \(B\) 时即封包。细包共 \(O(A)\) 个，每包负载至多为 \(2B\)。在冻结聚落，把各列表拼接起来。记 \(P=2^{64\lceil\log_2(2\Delta)\rceil}\) 及 \(K_1=\Delta P\)。由颜色的两两独立性，每个聚落中粗、细两类出向宽度之差的绝对值的期望至多为 \((B\sum_i a_i)^{1/2}\)，其中 \(a_i\) 是其第 \(i\) 个位置实例的出向核心宽度。求和可知，有一种着色使总不平衡量至多为 \(\sqrt{BQX}\)。

对每个源聚落，按目的包把粗出向单位分箱。由核心边两端的长度比可知，箱数至多为 \(K_1\)。把每个细包的所有出向单位视作一件物品，将整件物品匹配至任何一个剩余容量足够的箱中。首次匹配失败时，每个箱的剩余容量都小于 \(2B\)，故两侧总量中较小者的未匹配部分至多为 \(2BK_1\)。因此，按单侧计的匹配宽度满足 \[
 I\ge \tfrac12X-\tfrac12\sqrt{BQX}-2BK_1Q.
\] 聚落数上界~\eqref{eq:core-colonies} 与取 \(c_*\ge A\) 时的余量条件给出 \(BQ=O((X+Z+\sigma)/A)\)。取 \(A\) 为 \(1+K_1\) 的充分大的二的幂倍数，必要时取整。于是 \(I\ge X/4-O(Z+\sigma)\)，隐含常数为绝对常数。

在每个箱内，将已匹配的细单位与互异的粗单位配对。应用引理~\ref{lem:powered-gadget}，取 \(b_*=b\) 和 \(z=u/b_*\)；由细分假设后者为整数。在每一对中，粗线充当 \(AB\)，其源为 \(A\)，细线充当 \(CD\)，其源为 \(C\)。两个源属于同一聚落。由于一条线始终留在一个块内，且整个块颜色相同，两个目的端 \(B,D\) 分别属于粗包和细包。为每一对引入新的辅助方 \(M,N'\)。遍及所有配对，两条需求连线的宽度总和为 \(2I\)，每个乘积型求和项中替代线的总宽度至多为 \(2\lambda I\)。

在其中一个求和项内，对所有配对按行收缩后保留的宽度求和，再对按列收缩的情形作同样计算。每条替代线在两种布局之一中成为内部线，因此这两个总量之和至多为 \(2\lambda I\)。在所有配对中统一执行代价较小的那种收缩。进一步认同顶点只会使保留的宽度减小。未匹配连线的贡献至多为 \(X-2I\)，故 \[
 X'\le X-(2-\lambda)I\le\rho X+C(Z+\sigma),\qquad
 \rho=1-(2-\lambda)/4<1 .
\]

以整件为单位匹配每个细包，可以同时控制负载和度数的增长。被合并的旧包构成星形图。细包作为源时，在行布局中与自身所在聚落的粗包合并，在列布局中与其唯一选定的粗目的包合并。作为目的时，它仅与一个辅助方合并。因此每个粗中心至多吸收 \(O((1+K_1)A)\) 个细包。每个辅助方都并入一个需求端点，可以将其新增关联宽度按匹配宽度的常数倍计入该端点。可用一个二的幂次负载增长因子 \(g\)，同时控制所得核心负载和冻结列表的缩放负载，且该因子在 \(\Delta\) 的多项式以内。新单位为 \(u/b_*\)。

在每个可能的二进区间中预留包槽位。可能的核心接触只发生在同一聚落的包之间，或相交且长度比至多为 \(P\) 的相邻区间的包之间。对一个给定区间，在每个邻点有 \(O(P+\log(2P))\) 个可能区间。上述星形图及四个匹配端点在这张图中的直径都有界。取图的一个固定次幂，就得到度数至多为 \(C_2\Delta^{C_2}\) 的公共核心支撑图。系线只需经过一处具有自身长度比界的旧系线邻接，再经过有界多个核心邻接。因此全图最大度数至多为 \(C_2\nabla^{C_2}\)，其中绝对常数 \(C_2\) 必要时可增大。每个位置的候选区间少于 \(2r\) 个。这样便得到了所称的位置数界，并使各系线仍接至各自的冻结张量。

\subsection{重置核心度数}\label{subsec:reset} 以 \(H,N_b,\Delta,\nabla\) 表示重置前一刻的参数，并令 \(P=2^{64\lceil\log_2(2\Delta)\rceil}\)、\(K_1=\Delta P\)。重置用一个度数为三的宿主图替换核心。其主要开销是每个核心基区中的一个置换网络。固定一个规模为 \(c_0Yw_0\) 的三价网络，使得通过固定其标签，能够实现 \(Y\) 行（每行宽 \(w_0\)）的所有置换，其中 \(4K_1\le Y\) 且 \(c_0\ge1\)。称这一网络为\emph{转子}。

图~\ref{fig:reset} 给出两种路径。先根据平均负载，把每个位置的块划为称作\emph{基区}的区间，并在每对相邻核心基区之间预留容量。混排各块使大多数基区能符合这些预留容量。两个这样的基区之间的线，直接经过两个专用转子输出端之间的连接。与例外基区关联的每条核心线则沿二叉树走较长的路径。要将这些额外跳数计入预算，必须先控制例外基区的总负载。

\begin{figure}[htb]\centering \begin{tikzpicture}[x=1cm,y=1cm,
                        font=\footnotesize,term/.style={circle,draw,fill=white,inner sep=2pt}]
% pair of good bases
\foreach \d/\off in {1/-5,-1/5}{
  \begin{scope}[xshift=\off cm,xscale=\d]
    \draw[draw=black!60,rounded corners=6pt](-1.45,-1.5) rectangle (3.36,1.5);
    \node[font=\sffamily\small] at (1,1.8){良好基区};
    \node at (-.05,1.18){输入};
    \node at (2.6,1.18){输出};
    \foreach \y in {.75,0,-.75}{
      \draw (0,\y) -- (2.6,\y);
      \draw (0,\y) -- (-.65,\y) --(-.87,\y+.21);
      \draw (-.65,\y)--(-.87,\y-.21);
    }
    \draw[sea, very thick] (0,0)--(.8,0) (1.8,.75)--(2.6,.75);
    \draw[draw=ink,rounded corners,fill=ink!8] (.8,-1.26) rectangle (1.8,1.12);
    \draw[sea,very thick] (.8,0) -- (1.8,.75);
    \node[font=\scriptsize] at (1.3,-.59){转子};
    \foreach \y in {.75,0,-.75}{
       \node[term] at (0,\y){};
       \node[term] at (2.6,\y){};
    }
  \end{scope}}
\draw[sea,very thick] (-2.29,.75) -- node[below,text=black]{预留连线} (2.29,.75);

% link from bottom input on left to a terminal in a bad base
\draw[draw=black!60,rounded corners=5pt] (-2.15,-3.94) rectangle (0.2,-2.23);
\node[anchor=west,align=center] at (.42,-3.1) {不良基区的叶};
\node at (-.61,-2.7) {输入};
\draw (-1.41,-3.26) --(-1.68,-2.9);
\draw[rust,thick,densely dashed] (-5.11,-.75)--(-5.65,-.75)--(-5.87,-.96)
  -- (-5.87,-3.6)-- (-1.68,-3.6) -- (-1.41,-3.26)--(-.6,-3.26);
\node[term] at (-.6,-3.26){};
\node[align=center] at (-4,-2.91) {经接触树};
\end{tikzpicture} \caption{重置时的布线。在良好基区之间，连线的每个端点都通过一个转子抵达专用输出端（实线）。与不良基区的接触通过二叉接触树连接输入端（虚线）。旧张量及其系线留在输入端。图中只画出部分端口。}\label{fig:reset} \end{figure}

\paragraph{基区与容量。} 取二的幂 $F\ge Y^2w_0$，本步骤要求 $c_*\ge F$。对各位置在 \(N_b\) 个块上的平均负载 \(\mu_v\)，令 \[
 \nu_v=\max\left\{ FB/r,\ 
      \max_{\pi:y\to v} \mu_y\,2^{-\sum_{e\in\pi}s_e/2}\right\},
\] 允许长度为零的路径。与计数聚落时一样，从每个位置出发的各路径的衰减权重之和受一个绝对常数控制。因此 \(\sum_v\nu_v\le HFB/r+C\sum_v\mu_v\)。在位置 \(v\) 处分成等长基区，长度取满足 \(s_v\le FB/\nu_v\) 的最大二的幂，介于 \(F\) 与 \(r\) 之间。各次块指标的划分均与每 $r$ 个连续指标一组的分组相容。相邻基区的深度之差至多为 \(s_e/2\)。基区总数为 \(O((\mathcal M+\sigma)/(FB))\)。对相邻核心位置 \(v,w\) 上每对相交的基区，预留容量 \[
 \overline y_{vw}\min(s_v,s_w)+FB/K_1 ,
\] 其中 \(\overline y_{vw}\) 是按块平均的该接触的核心宽度。在任一基区，第一项之和至多为 \(s_v\mu_v\le FB\)，而预留项至多有 \(K_1\) 个，故该基区的总预留容量至多为 \(2FB\)。在置换各块之前，基区长度及容量已由均值固定。因而对每一组，给定一对基区位置的输出端将占用相同的宿主槽位。

\paragraph{混排。} 用同一个可逆二元矩阵置换所有块指标。在每个基区内，标记负载超过 \(2FB\) 的二进子区间（可能包括基区本身），按相同的边距离传播。叶区间负载至多为 \(2FB\)，长度至少为 \(F\)，并满足先前的长度比界（比较基区深度可知，标记区域上方的基区分裂已满足传播要求，再按聚落装包时的方法传播）。若基区被分裂，或与之关联的核心预留容量溢出，则称其为不良基区。对任意 \(\delta_*>0\)，按 \(\Delta,Y,\delta_*\) 选足够大的 \(F\)，即存在一个置换使 \begin{equation}
 \sum_{\text{不良基区}}\mathrm{load}\le\delta_*(\mathcal M+\sigma).
 \label{eq:core-badbases}
\end{equation} 为了使证明对全图最大度数一致，计算一个过载区间的传播范围。对有 \(n_v=N_b/s_v\) 个基区的某个给定位置，将原负载除以 \(FB\) 归一化，得到 \(x_i\le1/F\)。在低于基区 \(j\) 层的深度，设 \(a_J\) 为区间的归一化负载，\(\overline a\) 为其均值，它至多为 \(2^{-j}\)。量 \(V_j=\sum_J(a_J-\overline a)^2\) 和 \(U_j=\sum_J a_J^2-\sum_i x_i^2\) 各自均为过载区间数及其归一化负载总和的一半的上界。在随机可逆二元矩阵作用下，不同指标发生前缀碰撞的概率至多为区间数的倒数。因此 \[
 \mathbb E V_j\le\sum_i x_i^2\le n_v/F,\qquad
 \mathbb E U_j\le n_v2^{-j}.
\] 由于相邻基区的深度差至多为边距离的一半，深度 \(j\) 的标记只有沿满足 \(\sum s_e\le 2j\) 的路径才能迫使某个基区分裂。对该源标记，由二进区间的嵌套性，每个可达位置至多有一个基区受影响。因此由~\eqref{eq:core-colonies} 所用的路径计数，受影响数目至多为 \(2\cdot 2^{j/8}\)。在深度 \(j\) 处，若 \(1/F\le2^{-j}\) 则用 \(V_j\)，否则用 \(U_j\)。给这一估计量赋权 \(4\cdot 2^{\lceil j/8\rceil}\)，在 \(j=\log_2 F\) 处分开求和。期望计费量至多为 \(n_v\sum_{j\ge0}2^{j/8}\min\{F^{-1},2^{-j}\}=O(n_vF^{-7/8})\) 的一个常数倍。负载超过 \(2FB\) 的基区直接由深度零的估计支付。对核心容量溢出，在基区相交的深度，对接触宽度除以 \(FB\) 的结果应用方差估计量。预留容量比均值高出 \(FB/K_1\)，保证每次溢出至少贡献 \(1/K_1^2\)。将无权方差的期望对核心边求和，成本至多为 \(X/(F^2B)\)。每次溢出至多影响两个基区。其中负载超过 \(2FB\) 的已在深度零处计费，其余每个受影响基区的成本至多为 \(2FB\)。令 \(\Phi\) 为加权区间估计量之和再加上各接触方差之和的 \(K_1^2\) 倍。以上给出了不良基区总负载的确定性上界 $O(FB\Phi)$，并证明 \[
 FB\,\mathbb E\Phi = O(F^{-7/8}(\mathcal M+\sigma)+ K_1^2X/F).
\] 取估计量不超过其均值的一个矩阵，即得 \eqref{eq:core-badbases}。第~\ref{sec:assembly} 节通过计算碰撞得分找出这样的置换。

\paragraph{布线路径。} 混排后，将连续的 \(r\) 个块分成一组。给每个核心叶区间配置 \(Y\) 个输入端，把它的各旧张量逐个完整分配给当前负载最小的端口，使每个端口接收的负载至多为 \(2FB/Y+B\le 3FB/Y\)。所有系线都留在分配到的端口上，其另一端仍在拼接后的冻结列表中。在每个良好核心基区，转子将输入端连接至 \(Y\) 个输出端。对一项容量为 \(c\) 的预留，分配 \(\lceil cY/(8FB)\rceil\) 对专用输出端，每对均以直接连线相接。这样每个基区至多需要 \(Y/4+K_1\le Y\) 个端口。将两个良好基区之间一次接触的所有单位尽可能均匀地分到相应各对端口，输出负载至多为 \(8FB/Y+u\le9FB/Y\)。

在一个良好基区，记 \(m_{ij}\) 为从输入 \(i\) 路由到输出 \(j\) 的单位数。矩阵 \((\lceil m_{ij}/w_0\rceil)\) 的各行和及各列和均至多为 \[
 9FB/(Yuw_0)+Y\le10FB/(Yuw_0).
\] 通过添边把所有行和、列和都补到其中的最大值。由 Hall 定理，所得正则二部多重图可分解成完美匹配。每个匹配配置转子的一层，其中比特边取张量幂，使宽度达到 \(u\)。把未用的输入及其对应输出固定为零。在转子的每个位置，取各层表的张量积。这只用到通过固定转子标签获得的规定相等张量。在转子各行，平衡二叉树汇集其 \(w_0\) 个端口，每个单位经过 \(O(\log(2w_0))\) 跳。计入直接连线后，每个良好基区的开销为 \(O((c_0+\log(2w_0))FB)\)。

对一端不良的核心接触，使用接在相应输入端上的平衡二叉接触树。可能相邻的叶区间长度比至多为 \(P\)，每叶有 \(Y\) 个端口，故深度为 \(O(\log(2\Delta Y))\)。利用相等张量沿树传送有序列表。两端已分配的列表唯一确定所传比特的内部赋值。在 \eqref{eq:core-badbases} 中取充分小的 \(\delta_*\)，即足以支付该布线深度。因此，对某个绝对常数 \(A_0\)，有 \[
 X'\le A_0(1+c_0+\log(2w_0))(X+Z+\sigma).
\] 这里计入了转子、汇集树、直接连接以及接触树中的每一条核心连线。单位仍为 \(u\)。由端口与层的界，每个输入端、输出端、树顶点及转子顶点的负载都是 \(O(FB)\)。每个冻结叶区间中，缩放后的列表负载至多为 \(2FB\)。核心度数至多为三。转子顶点有此度数。在核心内，一个输入端至多与一个转子行（通过一条接往行树的边）及一棵接触树相连，一个输出端至多与一个转子行及一条直接连线相接。每棵二叉树度数至多为三。对系线，公共宿主图允许连接到长度比符合规定且可能相交的所有核心叶区间上的每个输入端。所得全图最大度数以 \(2\nabla\) 的某个固定次幂为界，其中的常数依赖于此次重置。在所有候选区间槽位上预留这些设施至多需要 \(c r H\) 个位置。

\subsection{支付重置的开销} 用 \[
 \kappa=A_0(1+c_0+\log(2w_0))
\] 表示重置的宽度乘数。取 \(t\) 为满足 \(\kappa\rho^t\le1/2\) 的最小正整数。反复使用子轮估计可知，在重置之前，核心宽度至多为 \(\rho^tX+C(Z+\sigma)/(1-\rho)\)。重置后，输入宽度 \(X\) 的系数至多为二分之一。其余开销为 \(Z+\sigma\) 的固定倍数，条件是对经过 \(t\) 个子轮后达到的度数，能选到满足 \(Y\ge4K_1\) 的转子。

起始核心度数有界时，递推 \(\Delta'\le C_2\Delta^{C_2}\) 给出 \[
 \log_2(4K_1)\le C^{t+1}
      \le C''(2+c_0+\log(2w_0))^{C'} .
\] 这里 \(C,C',C''\) 仅依赖于输入核心度数和各子轮常数。对 \(Y=2^y\)，由定理~\ref{thm:sort}，对每个 \(\zeta>0\) 可得 \(w_0=\lceil y^{a_\zeta}\rceil\) 及 \(c_0=O_\zeta(y^\zeta)\)。固定 \(\zeta<1/(2C')\)，再取充分大的 \(y\)。上面给出的多项式界至多为 \(y\)，故该转子可处理所需的核心度数。

子轮数 \(t\)、转子、\(F\) 及 \(\delta_*\) 仅取决于输入核心度数的上界和条件性的排序构造。这些参数控制最大负载的增长。全图最大度数决定系线传播的范围和要预留的宿主位置数。经过全部子轮，\(u\) 除以 \(b_*^t\)；各子轮中的多项式度数界再与重置合并，使 \(\log(2\nabla)\) 至多乘以一个固定因子。

引理~\ref{lem:core-width} 得证。


\section{一轮之后的元数预算}\label{sec:full}

集中操作所产生的冻结站点的总元数为 $O(L/\sqrt D)$。核心的最大元数很小，但所用的比特仍可能超过 $L$。反复执行各拍可缩减其内部宽度，直至总计费量落入较小的预算 $L/G$ 内。为重新执行集中操作，还要控制计费量的平方和，并使整个宿主图的度数有绝对常数上界。

\subsection{从集中到有界度数} 在引理~\ref{lem:local} 中取元数指数 $\ell=10$。该引理给出的核心度数绝对上界确定了引理~\ref{lem:core-width} 中一拍的参数。选取足够小的 $\epsilon\in(0,1/100]$，使得经过 \[
 J=\lceil2\epsilon\log D\rceil
\] 拍后，最大负载、线束单位的倒数、以及全图最大度数的对数这三个量各自的逐拍增长因子之积，在 $D$ 充分大时都不超过 $D^{0.03}$。每拍的这三个因子均与 $D$ 无关，所以可以这样选取。下文根据 $k,\delta,\epsilon$ 的取值，将 $D$ 增大到所需程度。

要求 $ed\le L/D$。在本轮的每个分组步骤中，还要求 \begin{equation}\label{eq:round-slack}
  H N_b T/r\le c_D L
\end{equation} 其中 $c_D>0$ 充分小，宿主图的数量参数取当前值，$L,T$ 取本轮开始时的值。下文将选取 $S,r$，使得以上条件在每轮均成立。

\paragraph{使元数的平方集中。} 将 $r$ 个块分组。在宿主度数有绝对上界时，对计费负载应用二进划分估计 \eqref{eq:core-colonies}，取最大负载 $DT/2$、容量 $DT$。由 \eqref{eq:round-slack}，遍及所有位置和组的非空区间共 $O(L/(DT))$ 个。把每个区间视为一个顶点，对所得图应用引理~\ref{lem:local}。将一切候选区间都列为位置。同一个旧位置与其旧邻点上的两个区间，只有在相交且长度比满足二进划分的绝对界时才允许相连。这包含了所有实际接触。二进区间的嵌套性使新宿主图的简单图度数以绝对常数为界。

在每个非空区间内，保留所有计费元数大于 $T$ 的站点。把其他站点依次装包，总计费元数一旦达到 $T$ 就封包。最后一包可以较小，每包元数至多为 $2T$。包的表等于其成员的表之积，其中所有内部指标都已收缩。每个区间由此得到 $O(D)$ 个站点。若区间有外部比特，则按原顺序移至另加的一个恒等站点，并将复制出的端口接回原先的关联位置。指定这些外部恒等站点为不可移动站点。它们的总数至多为非空区间数。在每个区间内，旧站点和各包的总元数至多为 $DT$，新恒等站点增加的元数至多为 $2DT$。原有大站点对平方和的贡献至多为 $LT$。每个新包的元数至多为 $2T$，因此各包的总贡献至多为 $2TL$。若区间 $i$ 的外部宽度为 $f_i$，那么新恒等站点的贡献至多为 \[
 4\sum_i f_i^2\le 4DT\,ed\le 4LT .
\] 取 $W=C_1L$，其中 $C_1$ 为充分大的绝对常数，就恰好满足引理~\ref{lem:local} 对占用数、不可移动站点及平方和的要求。具体地说，每个非空区间至多有 $D$ 个保留下来的大站点、$D+1$ 个包和一个外部恒等站点，符合该引理给出的上界 $100D$。

利用引理~\ref{lem:local} 展开。用 $Z$ 表示冻结站点的总计费元数，用 $X$ 表示核心站点间内部连线的总宽度，每条线只计一次。每个求和项中，冻结站点的计费量总和为 $Z=O(L/\sqrt D)$，平方和为 $O(LT)$，最大值为 $O(DT)$。每个核心站点的元数至多为 $D^{-10}T$，总元数为 $O(L(\log D)^{O(1)})$。这些估计中的常数均与 $D$ 无关。所有正宽度的内部线束均至少有宽度 $T/D^{O(1)}$。

为实行装包操作，取小于内部宽度保证下限的整数单位 $u$，使其仍满足 $u\ge T/D^{O(1)}$，且能被 $J$ 拍中规定的各细分因子整除。固定 $D$ 后，只要 $T$ 足够大，向下取这些因子之积的倍数即可。用固定为零的比特将每个正宽度内部线束向上填充到 $u$ 的倍数。元数及其平方分别至多增至原先的两倍和四倍。每个冻结站点的计费量升为原计费量与填充后元数的较大者，增长倍数至多为二，故总量 $Z$ 和平方和预算仍满足上述界。取覆盖所有核心站点、量级为 $D^{-10}T$ 的二的幂 $B$ 作为负载上界。仅在装包时，把冻结站点的计费量乘以一个量级为 $B/(DT)$ 的固定二进因子，该因子足够小，使缩放后的每一项至多为 $B$。实际宽度及冻结站点的平方和界都保持不变。将 $u$ 取得足够小，使 $u\le B$ 也成立。

\paragraph{消去核心连线。} 将引理~\ref{lem:core-width} 应用 $J$ 次，取余量 $\sigma=L/\sqrt D$。整个初始宿主图的度数在 $D$ 的多项式以内。中间各步的负载上界都是 $T$ 的固定倍数，倍数只依赖于 $D$，故每次由 \eqref{eq:round-slack} 得到所需余量。在引理~\ref{lem:core-width} 的递推式中，固定误差 $C(Z+\sigma)$ 乘上一个和小于二的几何级数。因此 \[
 X_{\mathrm{end}}\le 2^{-J}X_{\mathrm{start}}+O(Z+\sigma)
      = O\bigl(LD^{-2\epsilon}(\log D)^{O(1)}\bigr).
\] 冻结站点保留其计费量，不过现在一个位置实例可以容纳这些站点的一个列表。每个列表有 $D^{O(1)}$ 个成员。事实上，缩放后的总负载受最终负载界控制，而计费量的正下限给出每个成员缩放后负载的下界。每个核心站点的元数至多为 $D^{-9}T$。最终线束单位仍至少为 $T/D^{O(1)}$，核心的度数至多为三。当 $D$ 足够大时，全图最大度数 $\nabla$ 满足 $\log(2\nabla)\le D^{0.04}$。最后一项估计用到了度数对数的初始上界 $O(\log D)$。

\paragraph{恢复有界度数。} 单个站点的\emph{实际}邻点可能很少，但在公共宿主图中仍可有很多\emph{可能}的邻点。在冻结位置为其列表中每个成员保留一个槽位。冻结站点的接触全都是通向核心的系线。分三步为它们布线（图~\ref{fig:tether-trees}）。先把活动的系线单位依次分到独立出口。由最大宽度和最小单位可知，出口共有 $D^{O(1)}$ 个。再从各出口通过一棵二叉树寻址至任一可能的核心端点。在每个核心端点也设置一棵二叉树，对可能接入的出口寻址。连接两侧相应的叶节点。内部各分叉处的表，是其活动子节点的字符串拼接所得串上的恒等映射。不活动的分支无需站点。每个块按照自己的系线分配来使用这些树，而树和槽位本身在公共宿主图上固定不变。第一级树的深度为 $O(\log D)$，其余树的深度为 $O(\log(2\nabla)+\log D)=O(D^{0.04})$。宿主图的规模仅增加一个依赖于 $D$ 的因子。核心端点保留至多三个核心接触，再接一条树边，因此全图最大度数至多为四。

对线性预算，每个系线比特沿一条相应长度的路径传送，因而额外总元数为 $O(ZD^{0.04})$。对平方和预算，第一棵分线树受冻结站点元数控制。该树的每一层，总宽度至多为对应冻结站点的元数。因此该层的元数平方和至多为此元数平方的四倍。先对冻结站点求和，再对层数求和，得到 $O(LT\log D)$。其余每个新增顶点所承载的宽度至多是一个单位的宽度，或一个核心站点的全部系线比特数，因此元数至多为 $2D^{-9}T$。它对平方和的贡献至多为其元数的 $2D^{-9}T$ 倍。原有核心站点的总元数至多为 $2X_{\mathrm{end}}+Z$，并且同样满足这一估计。

\begin{figure}[tb]\centering \begin{tikzpicture}[x=1cm,y=.82cm,font=\footnotesize,
  live/.style={draw=sea,line width=1.3pt},
  possible/.style={draw=ink!55,densely dashed},
  branch/.style={circle,draw=ink,fill=white,inner sep=1.35pt},
  station/.style={rounded corners=2pt,draw=ink,inner sep=3pt}]
  \node[station,fill=rust!9] (f) at (0,0) {冻结};
  \node[branch] (split) at (1.2,0) {};
  \node[branch] (a) at (2.9,.78) {};
  \node[branch] (b) at (2.9,-.78) {};
  \node[branch] (c) at (6.45,.78) {};
  \node[branch] (d) at (6.45,-.78) {};
  \node[station,fill=ink!7] (c1) at (7.8,.78) {核心};
  \node[station,fill=ink!7] (c2) at (7.8,-.78) {核心};
  \draw[live,line width=2.3pt] (f)--(split);
  \draw[live] (split)--(a) (split)--(b);
  \draw[possible] (a)--(3.6,.43)--(5.76,-.43)--(d);
  \draw[possible] (b)--(3.6,-.43)--(5.76,.43)--(c);
  \draw[live] (a)--(3.6,1.12)--(5.76,1.12)--(c)--(c1);
  \draw[live] (b)--(3.6,-1.12)--(5.76,-1.12)--(d)--(c2);
  \foreach \x/\heading in {1.1/{拆成单位},3.9/{选择目的地},6.95/{在端点汇集}}{
      \node[font=\footnotesize] at (\x,1.75) {\heading};
  }
\end{tikzpicture} \caption{两个单位离开同一个冻结站点。实线路径在当前块中活动。虚线路径使两个出口在其他块中都能到达另一端点。左侧的分线树可能在 $O(\log D)$ 层内承载冻结站点的全部元数。中间每条活动分支承载一个单位。右侧的宽度受接收端核心站点的系线宽度控制。}\label{fig:tether-trees} \end{figure}

合并这些界，一轮结束时各计费元数 $z_v$ 满足 \begin{equation}\label{eq:round-costs}
 \begin{aligned}
 \sum_v z_v&=O\bigl(LD^{-2\epsilon}(\log D)^{O(1)}\bigr),\\
 \sum_v z_v^2&=O\bigl(LT(\log D)^{O(1)}\bigr),\qquad
                   \max_v z_v=O(DT).
 \end{aligned}
\end{equation} 对冻结站点，我们保留了引理~\ref{lem:local} 中的正下限。其余每个保留下来的站点都与一个正宽度内部线束相接，宽度至少为 $T/D^{O(1)}$。因此，\eqref{eq:round-costs} 中的所有计费量都可满足这一量级的下限。\eqref{eq:round-costs} 中的常数及 $\log D$ 的幂次均与 $D$ 无关。特别地，$ZD^{0.04}=O(LD^{-0.46})$ 这一项符合第一个估计。

令 $G=2^{\lfloor\epsilon\log D\rfloor}$，取二的幂 $q$ 使 $G^{1+\delta/4}\le q\le G^{1+\delta/3}$。当 $D$ 足够大时，这一区间包含一个二的幂。令下一轮从下列预算开始：\begin{equation}\label{eq:round-target}
                 L'=L/G,\qquad T'=qT .
\end{equation} 事实上，$D^{-2\epsilon}(\log D)^{O(1)}=o(1/G)$，而 $L'T'\ge LTG^{\delta/4}$ 足以吸收平方和预算中的对数损失。对充分大的 $D$，\eqref{eq:round-costs} 中的最大值至多为 $DT'/2$。最后固定充分小的二进常数 $\theta>0$，使 $\theta qT$ 不超过刚得到的正下限。选取 $G,q,D$ 时使用的估计均无需给 $\theta$ 指定下界，因为分组使用的是总负载上限，且在应用局部引理前已经把小站点重新装包。由此得证下一轮的不变量 \eqref{eq:round-invariants}。

\subsection{迭代及最终表的描述}\label{iter:subsection}

固定刚构造的一轮，包括 $D$、$G$、$q$ 和该轮的 $s$ 个分组步骤。为使总计费量节省一个 $h^a$ 的因子，取 \[
 j=\left\lceil\frac{a\log h}{\log G}\right\rceil,\qquad
 L_i=L_0G^{-i},\qquad T_i=T_0q^i \qquad(0\le i\le j).
\] 这里 $L_i,T_i$ 是完成 $i$ 轮之后的预算。还须选择 $r$ 和初始块数 $S$。较大的 $r$ 可降低只装了一部分的包的相对开销。较大的 $S$ 保证全部 $sj$ 步都有充足的组，还留下足够多的块，可在第~\ref{sec:assembly} 节中分摊公共宿主图的成本。

\paragraph{装包余量与外部宽度。} 经过一次分组，$N_b$ 除以 $r$，而 $H$ 至多增至原来的 $cr$ 倍。两者的乘积给出分配的位置实例数，其增长倍数至多为 $c$。宿主图的固定扩张可计入 $c$。因此在构造的每个阶段，都有 \begin{equation}
 H N_b\le C^{j+1}mS\log h , \label{iter:occurrences}
\end{equation} 其中 $C$ 仅依赖于单轮方案。因为 $L_0=C_0 mT_0 S\log(2+h)$，所以在以 $L_i,T_i$ 开始的轮次中，控制余量的比率满足 \[
   \frac{H N_b T_i}{rL_i}\le \frac{C_1C^{j+1}(qG)^i}{r}.
\] 取 $r$ 为不小于 $C_2(CqG)^{j+1}$ 的最小二的幂。由单轮方案固定 $C_2$，使得上述比率不超过 \eqref{eq:round-slack} 中的 $c_D$，并且 $r$ 满足每一步的下界。特别地，$r=h^{O_D(a)+o(1)}$，且与 $S$ 无关。外部宽度另有一个上界。每一轮都包含那 $ed$ 个外露比特，并且 $G^j<Gh^a$。因此对于 $i\le j$，有 \[
       \frac{ed}{L_i}=\frac{eG^i}{C_0m\log(2+h)}
           =O_D\!\left(\frac{h^{-a}}{\log h}\right).
\] 当 $h$ 足够大时，得到 $ed\le L_i/D$，恰为不可移动外部恒等站点预留的宽度。

\paragraph{副本与终止时的预算。} 任意固定 $M>0$，设 $R$ 是不小于 $h^M$ 的最小二的幂。取 $S=Rr^{sj}$，经过 $j$ 轮后恰剩 $R$ 个块。每个初始块含 $T_0$ 个副本，故 $d=T_0S$ 且 $\log d=O((\log h)^2)$。特别地，$d=n^{o(1)}$。现在 $h^a\le G^j<Gh^a$，$q\le G^{1+\delta/3}$，且 $T_0=\Theta(h^b)$。由 \eqref{eq:round-start} 和 \eqref{iter:occurrences} 得 \begin{equation}
 \begin{aligned}
 L_j&=O(md\log h\,h^{-a}),&
 T_j&=O\bigl(h^{b+(1+\delta/3)a}\bigr),\\
 HN_b&\le mS h^{O_D(a)+o(1)} .
 \end{aligned}
 \label{eq:round-terminal}
\end{equation} 位置实例数的界自始至终成立，且 $O_D(a)$ 中的常数与 $M,a,b,\beta$ 无关。由这些尺度也可知，整个过程中可以一直用同一个域。使用的尺度从 $T_0\ge4h^\beta$ 起，直至 $h$ 的一个固定次幂。选取扩张次数，使其为引理~\ref{lem:local} 中固定扩域次数的倍数，并且足够大，使域的基数在尺度 $T_j$ 上满足要求。其次数为 $O(\log h)$，故对每个 $i$ 也为 $O(\log T_i)$。

\paragraph{表的描述。} 每个位置实例携带其表的标签、各线束的宽度和这些线束在端口排序中的位置。每条关联宿主边可以用它在局部邻点列表中的位置命名，因此命名的成本可由 $D$ 控制。以 $D$ 为下标的常数可以依赖于单轮方案。每个位置实例的这些局部数据的总长度至多为 \begin{equation}
 C_3^{j+1}(\log h)^{O(1)}
   =h^{O_D(a)+o(1)} .
 \label{iter:labels}
\end{equation}

初始时，一个位置上的叠层有固定表标签，由其原表、叠层长度和各有序子列表的位置指定。这一描述的长度为 $(\log h)^{O(1)}$。在一轮之内，能够进入任何一个新标签的旧标签数，可由下述事实控制。\begin{itemize} \item 一个容量为 $DT_i$ 的聚落至多容纳 $D/\theta$ 个进入的活动站点，因为各站点的计费量至少为 $\theta T_i$。在引理~\ref{lem:local} 的局部展开中，每份输出描述使用有界个旧标签、离散选择及域元素。额外的描述长度为 $O_D(\log h)$。\item 在子轮或重置中，每个叶的负载上界，以及对同一新包有贡献的叶数上界，分别是 $T_i$ 和 $1$ 的固定倍数，倍数依赖于 $D$。每个活动核心至少贡献一个宽度为 $u\ge T_i/\operatorname{poly}(D)$ 的内部单位。单个冻结站点的计费负载即使经过固定缩放，仍有同类下界。标量站点一旦形成，随即取出。因此，每个有贡献的包都含有有界个旧标签和有界个单位。于是，单位配对列表、转子层列表以及一个端口处的有序子列表，其条目数都有界。每项至多需要 $O_D(\log h+\log r)$ 比特。这也描述了冻结站点处最终各棵树的布线。\end{itemize} 因此，若 $A$ 为当前标签长度的上界，则一步之后的上界为 $C_4A+C_4(\log h)^{C_5}$，其中包括在局部取出的标量因子的描述符。乘法常数计数的是旧标签，与扩域次数和输入字典标签的长度无关。迭代全部 $sj$ 步即证式~\eqref{iter:labels}。宿主图的全局地址只在控制整个位置时用到，无须附加在单个表的描述上。

\paragraph{总和所表示的张量。} 此时每一项由 $R$ 个保留下来的块以及在其演变过程中取出的各标量因子组成。重新分组时，外部列表与各自的站点一同置换，各乘积的每个表项都保持不变。局部张量恒等式把一个乘积换成一个和，而那 $ed$ 个已命名端口上的各个表项保持相同。着色、装包和布线在遇到平局时都按固定规则选择。唯有局部张量恒等式中的选择用来标记展开项。收缩最终的 $R$ 个块时，对最后的内部连线赋值求和。

一次这样的求值只需一部字典。一个键记录表的类型、阶段、描述和元数，并附上与表关联的比特。通过记录长度的字段或固定填充，得到唯一编码。输入字典和构造规则给每个有效键赋予域中的一个值。未用的键可任意赋值。特别地，已取出的标量的键没有关联比特。由最大元数不变量及式~\eqref{iter:labels}，最长键的长度为 \[
 O_D(T_j)+h^{O_D(a)+o(1)}
\] 推导出的表描述可能超过 $(\log h)^{O(1)}$ 比特。第~\ref{sec:assembly} 节把这段描述放进字典调用的\emph{实参}。安装的表的标签只需指定键的格式和长度。这样就分别计算了查表操作的元数，以及下一次宏操作要求输入的短标签。

